
\documentclass[12pt,a4paper]{ctexart}

\usepackage[margin=2.3cm]{geometry}
\usepackage{amsmath}
\usepackage{array}
\usepackage{booktabs}
\usepackage{caption}
\usepackage{enumitem}
\usepackage{etoolbox}
\usepackage{fancyhdr}
\usepackage{float}
\usepackage{graphicx}
\usepackage{hyperref}
\usepackage{longtable}
\usepackage{tabularx}
\usepackage[table]{xcolor}
\usepackage{tikz}
\usetikzlibrary{arrows.meta,positioning,shapes.geometric,fit,calc}

\hypersetup{hidelinks}
\urlstyle{same}
\setlist{nosep}
\captionsetup{font=small}
\renewcommand{\arraystretch}{1.25}
\setlength{\parindent}{2em}
\setlength{\parskip}{0.3em}
\setlength{\tabcolsep}{4pt}
\emergencystretch=3em
\sloppy
\makeatletter
\g@addto@macro{\UrlBreaks}{\do\/\do\-\do\_\do\.\do\:}
\makeatother
\AtBeginEnvironment{longtable}{\small\setlength{\tabcolsep}{3pt}}
\AtBeginEnvironment{tabularx}{\small\setlength{\tabcolsep}{3pt}}
\AtBeginEnvironment{tabular}{\small\setlength{\tabcolsep}{3pt}}

\newcommand{\studentname}{许奕}
\newcommand{\studentid}{2351441}
\newcommand{\teamname}{许奕、王天一、马源胜、周由}
\newcommand{\coursename}{软件工程管理与经济}
\newcommand{\projectname}{基于大模型的建筑能源智能管理与运维优化系统}
\newcommand{\docversion}{V2.0}
\newcommand{\docdate}{2026年5月31日}
\newcommand{\code}[1]{{\footnotesize\nolinkurl{#1}}}
\newcolumntype{Y}{>{\raggedright\arraybackslash}X}
\newcolumntype{C}[1]{>{\centering\arraybackslash}p{#1}}
\newcolumntype{L}[1]{>{\raggedright\arraybackslash}p{#1}}

\tikzset{
  box/.style={rectangle, rounded corners, draw=cyan!55!black, fill=cyan!8, thick, minimum height=8mm, align=center, inner sep=5pt},
  process/.style={rectangle, rounded corners, draw=teal!55!black, fill=teal!8, thick, minimum height=8mm, align=center, inner sep=5pt},
  data/.style={cylinder, shape border rotate=90, aspect=0.25, draw=orange!70!black, fill=orange!10, thick, minimum height=12mm, align=center, inner sep=5pt},
  risk/.style={diamond, draw=red!60!black, fill=red!8, thick, aspect=2, align=center, inner sep=2pt},
  arrow/.style={-{Stealth[length=2.4mm]}, thick, draw=gray!70!black},
  note/.style={rectangle, draw=gray!60, fill=gray!6, rounded corners, align=left, inner sep=6pt}
}

\pagestyle{fancy}
\fancyhf{}
\fancyhead[L]{\coursename}
\fancyhead[R]{\reporttitle}
\fancyfoot[C]{\thepage}
\setlength{\headheight}{15pt}

\title{
    \vspace{-2em}
    \heiti \zihao{2} \reporttitle \\
    \vspace{1em}
    \zihao{4} 课程期末项目正式交付文档
}
\author{
    \songti \zihao{4}
    项目名称：\projectname \\
    项目组成员：\teamname \\
    负责人：\studentname \quad 学号：\studentid
}
\date{\docdate}

\newcommand{\reporttitle}{最终提交说明}
\begin{document}

\maketitle
\thispagestyle{fancy}

\begin{abstract}
本文档为《\projectname》的最终提交说明，按照课程期末要求列出应提交的源码、SRS\&SDS、SEE\&SEM、测试与验收材料、演示材料和打包排除项。本文档用于最终检查，确保提交给教师的材料完整、正式、可运行且不泄露敏感信息。
\end{abstract}

\tableofcontents
\newpage

\section{教师要求对照}
\begin{longtable}{p{5cm}p{5.2cm}p{3.2cm}}
\caption{教师要求与项目材料对应关系}\\
\toprule
教师要求 & 对应材料 & 状态 \\
\midrule
\endfirsthead
\toprule
教师要求 & 对应材料 & 状态 \\
\midrule
\endhead
Complete financial analysis for course project & SEE 软件经济分析与评价 & 已完成 \\
Design, implement and deploy course project & 源码、README、用户手册与部署说明 & 已完成 \\
Project scope and planning & SEM 软件工程管理说明 & 已完成 \\
Software process monitor and control & SEM 软件工程管理说明 & 已完成 \\
Risk management & SEM 风险管理章节 & 已完成 \\
Cost of development and maintenance & SEE 成本估算与维护成本章节 & 已完成 \\
Pricing strategy & SEE 定价策略章节 & 已完成 \\
Funding and financing & SEE 资金与融资章节 & 已完成 \\
Financial analysis and evaluation & SEE 财务评价、ROI、敏感性分析 & 已完成 \\
SRS\&SDS & SRS 软件需求规格说明书、SDD/SDS 软件设计说明书 & 已完成 \\
SEE\&SEM & SEE 软件经济分析与评价、SEM 软件工程管理说明 & 已完成 \\
Source code developed and deployed & \code{source-code/}、README、启动脚本 & 已完成 \\
\bottomrule
\end{longtable}

\section{推荐提交目录}
\begin{verbatim}
期末提交材料/
|-- source-code/
|-- documents/
|   |-- latex/
|   |-- pdf-latex/
|   `-- markdown/
|-- demo/
`-- submission-readme.md
\end{verbatim}

\section{正式 PDF 清单}
正式 PDF 以 LaTeX 编译版为准：
\begin{enumerate}
    \item SRS 软件需求规格说明书。
    \item SDD/SDS 软件设计说明书。
    \item SEE 软件经济分析与评价。
    \item SEM 软件工程管理说明。
    \item 测试与验收报告。
    \item 用户手册与部署说明。
    \item 数据、接口与 MCP 说明。
    \item 最终提交说明。
\end{enumerate}

\section{源码提交要求}
源码应包含 \code{backend/}、\code{frontend/}、\code{data/}、\code{knowledge_base/}、\code{scripts/}、\code{README.md}、\code{.env.example} 和 \code{.gitignore}。源码应能够按 README 启动并运行完整演示。

\section{必须排除的内容}
\begin{longtable}{p{4cm}p{9.8cm}}
\caption{提交排除项}\\
\toprule
排除项 & 原因 \\
\midrule
\endfirsthead
\toprule
排除项 & 原因 \\
\midrule
\endhead
\code{.env} & 可能包含真实 API Key，严禁提交 \\
\code{frontend/node_modules/} & 依赖体积大，可通过 npm install 还原 \\
\code{frontend/dist/} & 构建产物可重新生成 \\
\code{data/runtime/} & 本地演示工单状态，不属于基线数据 \\
\code{.pytest_cache/} & 测试缓存，无提交价值 \\
个人无关实验文件 & 与期末项目无关，避免干扰评审 \\
真实 API Key 或截图 & 存在安全风险 \\
\bottomrule
\end{longtable}

\section{提交前检查}
\begin{enumerate}
    \item 运行 \code{scripts/check-project.ps1}，确认后端测试和前端构建通过。
    \item 检查最终目录中不存在 \code{.env} 和真实 API Key。
    \item 打开 LaTeX PDF，确认封面、目录、图表、页码和中文显示正常。
    \item 按用户手册启动系统，走一遍演示路径。
    \item 检查 PPT 和演示视频是否与当前系统一致。
    \item 将材料打包为一个 zip 文件并提交到 Canvas。
\end{enumerate}

\section{最终结论}
项目已经形成可运行源码、正式需求与设计文档、经济分析与工程管理文档、测试验收材料、用户部署说明和最终提交清单。软件侧已经具备最终提交条件，课程提交时只需补齐 PPT、演示视频并按本说明打包。

\end{document}
